\documentclass[10pt]{article}
\usepackage[letterpaper,landscape,margin=0.45in]{geometry}
\usepackage{array}
\usepackage{booktabs}
\usepackage{longtable}
\usepackage{xcolor}
\usepackage{hyperref}

\definecolor{slate}{HTML}{334155}
\definecolor{muted}{HTML}{64748B}
\definecolor{lightrow}{HTML}{F8FAFC}

\hypersetup{colorlinks=true, linkcolor=slate, urlcolor=slate}
\renewcommand{\arraystretch}{1.15}
\setlength{\tabcolsep}{2pt}
\setlength{\parindent}{0pt}
\setlength{\parskip}{4pt}
\newcolumntype{L}[1]{>{\raggedright\arraybackslash}p{#1}}
\emergencystretch=3em

\begin{document}

{\LARGE\bfseries Juplaya Trailer Power Lead Calculations}

{\small Planning worksheet for the TikZ power overview. This is not a final wiring schematic or code sign-off. Fuse class, interrupt rating, insulation type, conduit/bundle derating, crimp hardware, strain relief, and exact route lengths still need the measured install drawing.}

\section*{Inputs and Hot-Day Assumptions}

\begin{tabular}{@{}p{0.24\linewidth}p{0.70\linewidth}@{}}
\toprule
Item & Value used \\
\midrule
Primary sources & Project source files listed below the assumptions table; numeric values come from the current power design and local component spec notes. \\
Peak environment & Hot summer design ambient: 46 C / 115 F. Voltage-drop calculations use 75 C conductor resistance, so the assumed conductor rise for the drop calculation is \(\Delta T = 29\) C above ambient. \\
Copper temperature correction & \(R_T = R_{20}[1 + 0.00393(T-20)]\), with \(T=75\) C. \\
Voltage drop & \(V_{drop}=2LI R_T/1000\), where \(L\) is one-way length in ft and \(R_T\) is ohms per 1000 ft at 75 C. \\
PV module current & LG455N2W-E6: \(I_{mpp}=10.83\) A and \(I_{sc}=11.39\) A. Series strings keep the same current. \\
Charge-current convention & PV-string labels show string current. MPPT-to-battery labels show approximate 48 V charge current and controller output limit. \\
Wire-size convention & ``Minimum / plan'' means the smallest conductor selected for this planning model after ampacity and voltage-drop screening. It is not permission to ignore OCP, terminal temperature, bundling, conduit fill, or manufacturer instructions. \\
\bottomrule
\end{tabular}

{\footnotesize Source files: \texttt{docs/power.md}; LG455, Victron SmartSolar, Anker C1000/PS400, LiTime ComFlex, Dometic CFX3, and Orion reference notes under \texttt{docs/reference/}; Velit current range from \texttt{docs/research/aio-inverter-3panel-evidence-2026-06-05.md}.}

\section*{Peak Current Notes}

\begin{itemize}
\item Roof PV: \(3 \times 455\) W = 1365 W. Approximate 48 V charge current is 1365 W / 48 V = 28.4 A; at 57.6 V absorption it is 23.7 A. The SmartSolar 250/60-Tr output limit is 60 A.
\item Optional LG ground PV: \(2 \times 455\) W = 910 W. Approximate 48 V charge current is 19.0 A at 48 V and 15.8 A at 57.6 V. The SmartSolar 150/35 output limit is 35 A.
\item Velit 2000R: source notes carry about 5--14 A at 48 V, rated about 670 W. Use 14 A for branch voltage-drop sizing until the final manual/fuse value is confirmed.
\item 24 V bus: the Orion-Tr 48/24-16A is the hard rail limit. Fridge + normal lighting + USB fits; all floods, full awning, USB-C full load, optional C1000 top-up, and heater glow do not fit simultaneously.
\item Phase 2 MultiPlus is deferred. If AC charging is added later, combined trailer charging must still be capped at or below the 100 A ComFlex charge limit.
\end{itemize}

\section*{Lead Sizing Worksheet}

{\scriptsize
\begin{longtable}{@{}L{0.18\linewidth}L{0.12\linewidth}L{0.07\linewidth}L{0.08\linewidth}L{0.12\linewidth}L{0.10\linewidth}L{0.11\linewidth}L{0.13\linewidth}@{}}
\toprule
Lead & Current basis & Circuit V & Assumed one-way length & Required ampacity basis & Minimum / plan conductor & Worst-case drop at 75 C & Notes \\
\midrule
\endfirsthead
\toprule
Lead & Current basis & Circuit V & Assumed one-way length & Required ampacity basis & Minimum / plan conductor & Worst-case drop at 75 C & Notes \\
\midrule
\endhead
Roof 3S PV string to SmartSolar 250/60-Tr & 10.83 A MPP; 11.39 A Isc & 126.3 Vmp & 20 ft & PV conductor and OCP at or above 20 A module series-fuse class & 10 AWG / 2.59 mm & 0.53 V / 0.4\% & 3S raises voltage; current stays one-panel current. \\
Optional LG ground 2S string to SmartSolar 150/35 & 10.83 A MPP; 11.39 A Isc & 84.2 Vmp & 50 ft & PV conductor and OCP at or above 20 A module series-fuse class & 10 AWG / 2.59 mm & 1.32 V / 1.6\% & Portable length assumption; measure actual deployable lead. \\
SmartSolar 250/60-Tr to 48 V bus & 24--28 A expected; 60 A controller limit & 57.6 V charge & 4 ft & 60 A controller output & 6 AWG / 4.12 mm & 0.23 V / 0.4\% & Size OCP to controller and conductor, not average harvest. \\
SmartSolar 150/35 to 48 V bus & 16--19 A expected; 35 A controller limit & 57.6 V charge & 4 ft & 35 A controller output & 8 AWG / 3.26 mm & 0.21 V / 0.4\% & Optional source, but full-output conductor basis. \\
Battery main OCP / shunt / 48 V bus & 100 A continuous pack limit & 51.2 V nominal & 3 ft & 100 A charge/discharge limit, plus interrupt-rating gate & 2 AWG / 6.54 mm & 0.11 V / 0.2\% & Main OCP class and interrupt rating remain open. \\
48 V bus to Velit 2000R branch & 5--14 A operating range & 48 V nominal & 15 ft & At least 14 A plus final Velit OCP value & 12 AWG / 2.05 mm & 0.81 V / 1.7\% & Use final Velit manual/fuse value when available. \\
48 V bus to Orion 48/24 input & About 9.2 A at 384 W output and 87\% assumed efficiency & 48 V nominal & 4 ft & 20 A input breaker class from build notes & 12 AWG / 2.05 mm & 0.14 V / 0.3\% & Efficiency is assumed pending exact 48/24 datasheet row. \\
Orion 24 V output to Blue Sea 5026 & 16 A converter output limit & 24 V regulated & 3 ft & 16 A rail output & 12 AWG / 2.05 mm & 0.19 V / 0.8\% & Converter is the bus bottleneck. \\
Fridge branch & 4.6 A running & 24 V regulated & 18 ft & 10 A branch fuse; 3\% voltage-drop target & 14 AWG / 1.63 mm & 0.51 V / 2.1\% & Matches current build note if one-way run stays near 18 ft. \\
Longest lighting/flood trunk & 8 A planning trunk load & 24 V regulated & 35 ft & 5 A per branch; trunk must match aggregate branch plan & 10 AWG / 2.59 mm & 0.68 V / 2.8\% & Long exterior runs drive the size; local zones may be smaller. \\
USB-C PD branch & 5 A full USB-C load & 24 V regulated & 12 ft & 7.5 A branch fuse & 14 AWG / 1.63 mm & 0.37 V / 1.5\% & Full USB-C load is a short-duration worst case. \\
Optional C1000 top-up from 24 V bus & 10 A C1000 low-voltage input limit & 24 V regulated & 8 ft & 10 A manual/fused branch & 12 AWG / 2.05 mm & 0.31 V / 1.3\% & Discretionary; shed before it starves fridge/house loads. \\
Winter heater rough-in & 1--2 A run; up to 11 A glow estimate & 24 V regulated & 25 ft & 15 A rough-in branch & 10 AWG / 2.59 mm & 0.67 V / 2.8\% & Docs mention 12 AWG; 10 AWG is the voltage-drop planning pick at this assumed length. \\
48 V bus to Orion 48/12 input & 5.8 A at 240 W output and 87\% efficiency & 48 V nominal & 4 ft & Input OCP sized to converter and conductor & 14 AWG / 1.63 mm & 0.14 V / 0.3\% & 12 V receptacles stay cabinet-local only. \\
12 V cabinet receptacle bank & 20 A / 240 W converter output & 12.2 V regulated & 3 ft & 20 A converter output, branches fused individually & 12 AWG / 2.05 mm & 0.23 V / 1.9\% & Keep physically short; do not distribute a 12 V house rail. \\
Anker PS400 to C1000 & 8.33 A operating; C1000 high-range input allows 12.5 A & 48 V operating & 15 ft & Anker cable/input rating & 10 AWG / 2.59 mm & 0.30 V / 0.6\% & Separate portable island; never combine with LG strings. \\
\bottomrule
\end{longtable}
}

\section*{Load-Shedding Implications}

The 24 V rail cannot be sized by summing every optional load and pretending they are simultaneous. A literal worst case is fridge 4.6 A + interior lights 3 A + all exterior floods 5.25 A + awning strip 2.7 A + USB-C 5 A + C1000 top-up 10 A = 30.55 A before heater glow. That is almost twice the 16 A Orion limit. The operating rule is therefore part of the electrical design: C1000 top-up, all-floods mode, full USB-C load, and heater glow are mutually managed loads.

The roof and ground MPPT outputs are the opposite case: normal current is below controller capacity, but the conductor/OCP basis should follow controller output limits. Roof plus optional ground charging at expected full-array output is roughly 40--47 A depending on battery voltage. If both controllers are allowed to their hardware limits, they can contribute up to 95 A, still below the ComFlex 100 A charge cap, but leaving no room for a future AC charger.

\section*{Open Measurements Before Final Wiring}

\begin{itemize}
\item Replace every assumed one-way length with measured route length after the nose cabinet, roof gland, Velit station, fridge bay, switch panel, and exterior lighting paths are fixed.
\item Confirm conductor insulation temperature rating, strand class, terminal temperature limits, conduit or loom bundling, and whether any wire runs share hot cabinet air.
\item Choose final 48 V OCP hardware with adequate DC interrupt rating. The main battery OCP remains a Class-T-or-equivalent gate.
\item Bench-measure SmartScan and winter-heater glow current before making them part of a simultaneous-load case.
\item After shakedown, measure voltage at the load while hot and loaded; the calculated drops are only as good as the assumed lengths and terminations.
\end{itemize}

\end{document}
